\documentclass[conference]{IEEEtran}
\usepackage{booktabs,graphicx,amsmath,url,xcolor,array,enumitem,caption,balance}
\usepackage[hidelinks]{hyperref}
\setlist{nosep}
\newcommand{\code}[1]{\texttt{#1}}
\newcommand{\tool}{\textsc{RuleScope}}   % anonymized tool name for review (not the repository name)
% ISPASS 2027: 9 pages including figures/appendices, unlimited references, double-blind.
\title{Same Rewrite, Different Sign:\\Blaming and Guarding Hardware-Dependent Graph Rewrites in TorchInductor}
\author{\IEEEauthorblockN{Anonymous Authors}\IEEEauthorblockA{Paper \#\_\_\_}}

\begin{document}
\maketitle

\begin{abstract}
TorchInductor applies a few hundred graph-rewrite rules whenever they match, with no check that the rewrite is
faster on the target. We make every rule, flag and opt-in pass of TorchInductor individually switchable from
outside the compiler, calibrate a timing oracle on 480 null runs (false alarms 9\% for one comparison, 1.5\% with
two time-separated confirmations), and attribute a measured slowdown to the responsible rule or pair by delta
debugging over the switch set. On 24 models and six machine/precision targets (Apple M4 Pro, Intel Xeon, NVIDIA A40
in fp32 and bf16, L4, RTX 4090) eight of the sixteen most impactful switches have the opposite dominant sign on
different machines, while the structural change the compiler makes is identical; attribution recovers planted
culprits in 5--16 oracle calls and finds genuine pair interactions, and an uncalibrated oracle fabricates them.
Acting on the blame, learned per-rule guards remove the harm where it is frequent and shape-predictable (CPU matmul
decomposition from 261\% to 1\% lost per firing; MobileNet 29\% and EfficientNet 41\% faster) and, where a rule is
right most of the time, the measurements fix a default: \code{freezing}, off by default, is faster on 105--113 of 163
real networks on each GPU (median 7--10\%), a finding we filed upstream together with a Dynamo bug in the setting
that would remove its memory cost.
\end{abstract}

\section{Introduction}
% from report Sec 1, compressed to ~0.8 page: problem, thesis, contributions (4 bullets), results summary
\code{torch.compile} lowers a PyTorch model through TorchInductor, and a large part of what Inductor does is graph
rewriting: a few hundred rules each look for a pattern (a manual attention computation, a convolution followed by
batch normalisation, a matrix multiplication of a particular shape) and replace it with something the rule's author
expected to be faster. Rules fire whenever they match. Nothing measures whether the result is faster on the target,
and when a compiled model is slower than expected there is no way, for users or for the compiler's developers, to
find out which rule did it~\cite{rfc}. The only published per-optimization view of Inductor is an eight-knob
ablation on one GPU reported as a geometric mean~\cite{pt2}.

This paper makes the question answerable and reports what the answer looks like on six machine/precision targets.
Our thesis is that a graph rewrite fixes the \emph{shape} of the generated program, and whether that shape is faster
depends on the relative cost of library calls, generated kernels and memory traffic on the target, which no rule
written once can know. Three measured facts support it: half of the most impactful rules are confirmed to help on
one machine and hurt on another (\S\ref{sec:signs}); when they flip, the compiler has made the identical structural
change on both machines (Table~\ref{tab:mechanism}); and acting on the measurement pays, either through a learned
guard per rule where harm is frequent and predictable or through a per-machine default where it is not
(\S\ref{sec:guards}).

\noindent\textbf{Contributions.}
\begin{enumerate}[leftmargin=*]
\item \tool: per-rule switches for every pattern rule, configuration flag and opt-in pass of TorchInductor 2.14
  (269/396/398 switches on the Apple, x86 and CUDA builds), installed from outside the compiler by wrapping the hook it
  already calls before applying a rule, with fire and suppression counts (\S\ref{sec:switches}).
\item A timing oracle whose false-alarm rate was measured on 480 runs of unchanged configurations and reduced from
  9\% to 1.5\% by two time-separated confirmations, and delta-debugging attribution over switch sets that recovers
  planted culprits in 5--16 oracle calls, detects pairs, and declines to blame when nothing is slower
  (\S\ref{sec:attr}).
\item A study of 24 models on an Apple M4 Pro, an Intel Xeon, and NVIDIA A40, L4 and RTX~4090 GPUs in fp32 and bf16:
  the rule$\times$machine sign matrix, its mechanism, and the effect of precision on which rules fire at all
  (\S\ref{sec:signs}).
\item Two ways to act on blame, both evaluated on held-out models against a per-model oracle: learned per-rule guards,
  and per-machine defaults; the latter yields a concrete upstream finding about \code{freezing} backed by 163 real
  networks per GPU, and a Dynamo bug (\S\ref{sec:guards}).
\end{enumerate}


\section{Per-Rule Switches}
\label{sec:switches}
Inductor applies rewrites in pre-grad, joint-graph and post-grad passes; most are \code{PatternMatcherPass} instances
holding \code{PatternEntry} rules keyed by the operator they match, and a pass applies a rule only if
\code{entry.extra\_check(match)} returns true. \tool\ replaces every entry's \code{extra\_check} with a gate that refuses
the match when the rule's switch is off, counts a \emph{fire} when the original check accepts and the switch is on, and
counts a \emph{suppression} when the rule would have fired but was disabled. No PyTorch source is modified; Inductor's
own in-tree learned heuristics use the same hook from the inside~\cite{autoheuristic}. Rules register lazily, so
discovery compiles a tiny warm-up graph (attention, conv/batch-norm, split/cat, matmul, a dtype round trip) and then
walks the \code{fx\_passes} modules; direct calls to the initialisers register patterns twice. Configuration flags
(\code{freezing}, \code{layout\_optimization}, \code{max\_fusion\_size}, \ldots), opt-in passes enabled through the
fusion-option dictionaries, and master switches for whole pass groups complete the universe. A \emph{state} is the set
of switches that are on; because universes differ per platform (oneDNN rules exist only on x86, half-precision
attention variants only with CUDA) states are compared by their difference from the default. The FX-graph cache would
otherwise return a program compiled under another state, so a custom pass whose \code{uuid()} hashes the active state
keeps the cache key honest, and the cache is disabled while timing.

Most rules never match a given model: a \emph{verify} step flips each switch alone from the default, compiles, and
hashes the normalised generated code; for BERT about ten of 398 switches change the program. Everything downstream
works on that short list.


\section{A Calibrated Oracle and Delta-Debugging Attribution}
% TODO: cut from report/final/final.tex Section


\section{The Same Rewrite on Six Targets}
\label{sec:signs}
Fig.~\ref{fig:signs} and Table~\ref{tab:signmatrix} give, for the sixteen most impactful switches, on how many corpus
models having the switch ON is confirmed to hurt or help on each target. \textbf{Eight of sixteen have the opposite
dominant sign on different machines.} Layout optimisation hurts on the Mac (MobileNet 41\%, EfficientNet 68\%) and the
A40 in fp32 (ResNet-18 9\%) and helps on the Xeon (up to 90\%) and every bf16 GPU; the fusion-size limit helps
everywhere but the Xeon; matmul decomposition costs 135\% on the Mac and saves 95\% on the Xeon; freezing helps
wherever it does anything.

\begin{figure*}[t]\centering\includegraphics[width=0.62\textwidth]{figures/sign_matrix.pdf}
\caption{Rule$\times$machine sign matrix: colour is the balance of models helped vs.\ hurt by the switch being ON;
cells give helped/hurt counts (confirmed, $\geq3\%$).}\label{fig:signs}\end{figure*}
\input{tables/sign_matrix}

\textbf{Mechanism.} Joining each confirmed verdict with the code-level change the same switch makes
(Table~\ref{tab:mechanism}) shows the compiler making the same structural change on every machine.
\code{decompose\_mm} removes 96--106 external BLAS calls on both CPUs: 39\% slower on the Mac, whose BLAS handles
small matmuls well and whose generated loops do not, 53\% faster on the Xeon. Layout optimisation adds two kernels and
4--22\,MB of allocation on both CPUs because the batch-norm/ReLU epilogues stop fusing into the convolution: $-28\%$
on the Mac, $+14\%$ on the Xeon whose oneDNN convolutions prefer channels-last. The fusion limit removes 8--12 kernels
on the GPUs and is worth 7--12\% on 11--13 models per GPU. Relay noted in a sentence that pass effectiveness is
``model- and device-dependent''~\cite{relay}; this is the per-rule, six-target version of that sentence.
\input{tables/mechanism}

\textbf{Precision changes the rule set.} In fp32 no attention-fusion rule fired on any GPU: a GPU-only softmax
rewrite fires first and the fusion pattern no longer matches. In bf16 the half-precision variant of rule 28 fires on
BERT and is worth 7.5\% on the L4. A GPU study in fp32 would have concluded that attention fusion never fires.


\section{Acting on the Blame}
\label{sec:guards}
\textbf{Learned guards.} A guard is a small model in the same \code{extra\_check} hook that decides whether to let a
rule fire. Rows are (features at the firing, measured effect of blocking it), from synthetic sweeps (attention B/T/H/D,
conv C/R/K/L/depthwise, decode-MLP M/I/Hd; 1,516 rows on the CPUs) and from real networks (38 torchvision conv nets,
14 transformer architectures; 163 rows per switch per GPU, Table~\ref{tab:datasetsgpu}). Features are the shapes at
the match for pattern rules and a whole-model profile (conv work and its depthwise/1$\times$1 fractions, linear work
and how decode-shaped it is, attention shapes, parameters, batch, tokens) for model-level decisions. Trees,
gradient-boosted trees, logistic regression, a fine-tuned encoder and a commercial decision API are compared on grouped
splits that hold out whole shape slices or architectures; the score is run time lost per firing against
``always fire''. On the CPUs the guards lose 1.0--1.4\% per firing where always-fire loses 262\% (decode MLP) and
176\% (layout); end to end the layout guard makes MobileNet 29\% and EfficientNet 41\% faster on the Mac with no
regression, and the matmul guard with the decomposition pass enabled makes SmolLM decode 58\% (Mac) and GPT-2 decode
73\% (Xeon) faster (Table~\ref{tab:e2e-mac_mlp}). On the A40 in fp32 a layout guard trained on 60 real networks made
held-out ResNet-18 8\% and ResNet-50 4\% faster, recovering 97\% of the per-model oracle on the conv stack.
\input{tables/guard_cv_combined}
\input{tables/e2e_mac_mlp}
\input{tables/datasets_gpu}

\textbf{Where guards do not beat the default.} A guard trained on 1--3-layer synthetic stacks did not transfer to real
GPU networks until trained on real networks: the mechanism was fine, the distribution was not. On the bf16 GPUs the
layout rule hurts a minority of real networks (18--21 of 111 with an effect) and helps the majority (62--83), so no
learned guard beats ``always on'' in cross-validation there. Guards pay where harm is frequent and shape-predictable;
otherwise the measured answer is a per-machine default.

\textbf{A default: \code{freezing}.} It folds parameters into constants and is off by default. Across 163 real networks
on each of three GPUs it is confirmed faster on 105--113 and slower on at most one, median 7--10\%, p90 18--24\%, up
to 65\% (Table~\ref{tab:freezing}, Fig.~\ref{fig:freeze}); it helps 22 of 26 corpus models on the Xeon, up to 95\%, and
is a wash on the Mac, which has no oneDNN path. Verdicts and magnitudes are the same with published weights as with
random initialisation (Table~\ref{tab:pretrained}, 10--12 of 12 agree per machine). SmolLM-135M \code{generate} with a
KV cache is 8--22\% faster on an A5000 with identical output; ResNet-50 classification 4\% (GPU) and 11\% (Xeon). The
costs are compile time (+10--30\%) and peak memory (+45--51\% on decoder models). \code{freezing\_discard\_parameters}
removes the duplicate (ResNet-50 121$\to$75\,MB, BERT 280$\to$235) at no speed cost but crashes Dynamo on every
tied-embedding language model, which we reduced to ten lines on the CPU. Both were filed upstream (link withheld for
review).
\input{tables/freezing_zoo}
\begin{figure}[t]\centering\includegraphics[width=\linewidth]{figures/freezing_cdf.pdf}
\caption{Speed-up from \code{freezing} over 163 real networks per GPU (bf16).}\label{fig:freeze}\end{figure}
\input{tables/pretrained}


\section{Repeatability, Controls and Threats}
\label{sec:threats}
Same-machine verdicts repeat: 213 of 220 single-switch verdicts agree across sessions on the CPUs and 188 of 221 on
the A40; the disagreements are 2--3\% effects at the threshold and sub-millisecond toys. Across machines 42 of 47
disagree, which is the finding of \S\ref{sec:signs}. The control experiment of \S\ref{sec:attr} shows what an
uncalibrated oracle does: on the case whose slow state was within noise it reported a \code{batch\_tanh}$\times$
\code{decompose\_mm} interaction that does not exist, in eleven evaluations, where the calibrated oracle said
``not slower'' in two. PyTorch 2.13 registers 267 switches on the Mac against 269 for 2.14 (four rules added), and a ten-model
subset swept under 2.13 gives the same sign as 2.14 on every switch with data but the opt-in normalisation pass
(Table~\ref{tab:signmatrix}): the rule set drifts between releases, the hardware-dependence does not. Threats: inference
only; one compiler version for the full study; random weights for most of the corpus with 12 pretrained twins
agreeing; 163-row GPU datasets; a shared VM and rented pods, which is why every verdict needs two confirmations.


\section{Related Work}
Delta debugging is Zeller's~\cite{zeller}; applied over compiler passes it localises miscompiles with a Boolean oracle
(\code{opt-bisect}, BugLens, MultiConf), and performance bisection over commits is unreliable under noise~\cite{emse22},
which is why ours is calibrated. Pe\v{c}im\'uth et al.\ localise Graal performance problems by diffing recorded
optimisation decisions~\cite{mplr23}; we intervene and obtain a counterfactual culprit, closer in spirit to causal
profiling~\cite{coz}. Learned compiler heuristics span Milepost, MLGO~\cite{mlgo} and DeepTune~\cite{deeptune};
ML compilers learn cost models to drive schedule or rewrite search (Ansor~\cite{ansor}, TASO~\cite{taso},
Tensat~\cite{tensat}), and such models are known to be hardware-bound~\cite{tlp}; Inductor's own
\code{autoheuristic}~\cite{autoheuristic} learns one or two decisions per GPU. That optimisation benefit depends on the
target is old news at the flag level~\cite{pallister13,magni14}; none of these report per-rule sign flips in a production
ML compiler, and the only prior per-optimisation ablation of Inductor is Table~4 of the PyTorch~2 paper~\cite{pt2}: eight
knobs, one A100, geometric mean, no noise model, measured with freezing on. Benchmarking methodology follows
\cite{hoefler15,mytkowicz09,stabilizer,vmwarmup}.


\section{Conclusion}
A rewrite rule in a production ML compiler is a bet about hardware its author did not have. With per-rule switches, a
calibrated oracle and delta debugging the bet can be checked, the losing rule named, and the loss removed: by a guard
where it is predictable, by a default where it is not.

\balance

\bibliographystyle{IEEEtran}
\bibliography{refs}
\end{document}
